\syoumon{1}
$\sqrt{2} = \dfrac{q}{p}$ を満たす整数$p,q$ ($p\neq 0$) が存在すると仮定すると,
$2p^2 = q^2$ となる. 両辺が$2$で割れる回数を考えると, 左辺は奇数回, 右辺は偶数回であるので, 素因数分解の一意性により矛盾である.
よって, $\sqrt{2}$ は無理数である.

\syoumon{2}
$\log_{2}{10} = \dfrac{q}{p}$ を満たす整数$p,q$ ($p\neq 0$)が存在すると仮定する.
左辺は正なので, $p,q$ は正である.
このとき, $10^p = 2^q$ となるが, 左辺は$5$の倍数なのに右辺はそうではない. よって矛盾である.
\syoumon{3}
(1) とほとんど同様に示されるので, 省略.

\syoumon{4}
$t = \tan{1^{\circ}}$が有理数であると仮定する. このとき,
$\tan{2^{\circ}} = \dfrac{2t}{1-t^2}$ も有理数である.
同様にして,
$\tan{4^{\circ}}, \tan{8^{\circ}}, \dots, \tan{64^{\circ}}$ も有理数である.
$a = \tan{4^{\circ}}, b = \tan{64^{\circ}}$ とすると,
$$
\tan{60^{\circ}} = \sqrt{3} = \dfrac{b-a}{1+ab}
$$
で, 最右辺は有理数, $\sqrt{3}$は無理数であるので矛盾である.

\syoumon{5}
$\cos{1^{\circ}}$ が有理数であると仮定する.
$\cos{2\theta} = 2\cos^{2}{\theta} - 1$ より,
(4)と同様の手法で$\cos{(2^n)^{\circ}}$ ($n$は0以上の整数) は有理数である.

ここで, cosの積和公式は\footnote{sinをなるべく持ち出したくないので, cosで完結する公式を選んだ. }
$$
\cos{(A+B)} + \cos{(A-B)} = 2\cos{A}\cos{B}
$$
であった. いま非負整数$n$に対して $c_n = \cos{(n)^{\circ}}$ と置くと,
$$
c_{n+1} + c_{n-1} = 2c_n c_1
$$
である.
いま $c_1, c_2$ は有理数であり, 上の式より $c_{n}, c_{n-1}$ が有理数であれば$c_{n+1}$も有理数であるから, 帰納的に
$c_{30} = \dfrac{\sqrt{3}}{2}$ が有理数となるが, これは矛盾である.

\syoumon{6}
$x = 3^{\frac{1}{3}}p+\sqrt{2}q$ が有理数であると仮定する.
このとき$(x - \sqrt{2}q)^3 = 3p^3$ が有理数である. 左辺は
$$
x^3 - 3x^2\sqrt{2}q + 6xq^2 -2\sqrt{2}q^3
$$
であり, $x,q$ は有理数であるから, $\sqrt{2}$の係数部分である
$$
-3x^2\sqrt{2}q - 2\sqrt{2} q^3
$$
が有理数でなければならない.
これは
$$
-\sqrt{2}\cdot (3x^2q + 2q^3)
$$
となり, $q>0$ なので, $3x^2q + 2q^3$は正の有理数である.
よって上が有理数となることはなく, 矛盾であり, $x$は無理数である.

\noindent\textbf{Bonus.} 本問の主張において $p, q>0$は緩めることができます. 実際には,
$p, q$ は $(p,q)\neq (0,0)$ であれば無理数です.
練習として, 証明を書き換えてみてください.